\documentclass[11pt]{article}
\usepackage{amsmath,amsfonts}
\usepackage{graphicx}
\usepackage{xcolor}
\definecolor{gray}{rgb}{.75,.75,.75}
%\usepackage[landscape]{geometry}
%\usepackage{hyperref}
\usepackage[bookmarks=true, pdfborder={0 0 0}, urlbordercolor=gray]{hyperref}

\renewcommand{\thefootnote}{\fnsymbol{footnote}}

\addtolength{\topmargin}{-.6in}
\addtolength{\textheight}{.8in}
\addtolength{\oddsidemargin}{-.3in}
\addtolength{\textwidth}{.6in}
\pagestyle{empty}

\newcommand{\ds}{\displaystyle}
\newcommand{\ts}{\textstyle}

\newcommand{\R}{\mathbb{R}}
\newcommand{\tZ}{\tilde{\Z}}
\newcommand{\sgn}{\mathrm{sgn}}

\newcommand{\Sym}{\mathrm{Sym}}
\newcommand{\tu}{\textunderscore}

\renewcommand{\circle}[1]{{\bigcirc\hspace{-.75em}#1}}
\renewcommand{\epsilon}{\varepsilon}
\newcommand{\st}{\mbox{\footnotesize$\sqrt3$}}

\newcommand{\ul}{\underline}

\begin{document}

\footnote[0]{Notes by Peter Magyar \ {\tt magyarp@msu.edu}}
\vspace{-1em}

\centerline{\bf Math 132\hspace{1.5em} \hfill{\Large\bf More Uses for Integrals}\hfill \S4.4, 5.5}
\vspace{1em}

\noindent
{\bf Review.} 
The integral $\int_a^b f(x)\,dx$ has four levels of meaning.
\begin{itemize}
\item Numerical. 
To compute the {\it cumulative effect} of an influence function $y=f(x)$
over an interval $x\in [a,b]$, we define the {\it integral}.
Divide $x\in[a,b]$ into $n$ increments 
of width $\Delta x=\frac{b-a}{n}$, 
and choose sample points $x_1,\ldots, x_n$ inside each increment. 
Then the integral is the limit of Riemann sum approximations,
the total of the approximate contributions of each interval as $n\to\infty$, $\Delta x\to 0$:
$$
\int_a^b \!\!f(x) \,dx \ = \ 
\lim_{n\to\infty} \sum_{i=1}^n f(x_i)\, \Delta x
\ = \ 
\lim_{n\to\infty}
f(x_1)\Delta x +\cdots + f(x_n)\Delta x\,.
$$
\vspace{-1.3em}

\item Geometric. The integral is the {\it area} between the graph $y=f(x)$
and the interval $x\in[a,b]$, counting area above the $x$-axis as
positive, area below as negative.
The area  is the cumulative effect of the height function $f(x)$.

\item Physical. Take physical variables $y,z$ depending on an independent variable $x$, 
so that $y=f(x)$ and $z=F(x)$ for continuous functions $f$ and $F$.
We have two equivalent relationships between $y$ and $z$. 
First, $y$ is the rate of change of $z$:
$$y=\frac{dz}{dx} \qquad\text{and}\qquad f(x)=F'(x).$$
Second, the cumulative effect of $y$ is the change in $z$ over an interval $x\in[a,b]$.
For example, velocity $v$ is the rate of change of position $s$ per unit time, 
and the cumulative effect of velocity is change in position (distance traveled, displacement, $\Delta s$) 
over a time interval.
This is the {\it Second Fundamental Theorem of Calculus},
written in Leibnitz and Newton notation as:\footnote{
In Leibnitz notation, a function is denoted by its output variable, such as $z=F(x)$,
and its derivative function is $\frac{dz}{dx} = F'(x)$.
A particular output value of the function is denoted: $\left.\vphantom{\sum}z\right|_{x=a}=F(a)$;
and the change in the value over an interval $x\in[a,b]$
is denoted: $\left.\vphantom{\sum}z\right|_{x=a}^{x=b} = F(b)-F(a)$.
}
$$
\int_a^b\! y\ dx \ =\ \left.\vphantom{\sum}z\,\right|_{x=a}^{x=b}\ =\ \Delta z 
\qquad\text{and}\qquad
\int_a^b\! f(x)\,dx \ =\ F(b)-F(a)\,.
$$
Conversely, the rate of change of a cumulative effect up to a point
is equal to the strength of the effect at that point.
This is the {\it First Fundamental Theorem of Calculus}: 
the integral of $y=f(x)$ produces an antiderivative $z=I(x)$:
$$
I(x) = \int_a^x f(t)\,dt
\qquad\Longrightarrow\qquad
I'(x) \ =\ \frac{d}{dx} \int_a^x\! f(t)\,dt \ =\ f(x)\,.
$$
How to combine physical units of variables is clear in Leibnitz notation: $dy$ has the same units as $y$,
and the units factor out of derivatives and integrals. E.g.~for 
time $t$ in sec, velocity $v$ in meters/sec, we get acceleration and displacement as:
$$
\frac{dv}{dt} \,\frac{\text{ m/sec}}{\text{ sec}} = a \ \frac{\text{m}}{\text{sec}^2}
\quad,\quad \int_a^b \left(v(t)\frac{\text{m}}{\text{sec}}\right)\!(dt\text{ sec}) \,=\,  
\ \Delta s\, \frac{\text{m}}{\text{sec}}\text{sec}\ =\ \Delta s \text{ m}.
$$
\\[.8em]

\item Algebraic: Given $f(x)$, we can often reverse derivative formulas to find 
an antiderivative $F(x)$ with $F'(x)=f(x)$. 
Then the Second Fundamental Theorem immediately computes the integral $\int_a^b f(x)\,dx = F(b)-F(a)$. 
Later, we will develop techniques for finding $F(x)$, 
such as the Substitution Method which reverses the Chain Rule.

\end{itemize}

\noindent
{\bf Indefinite integral notation.} 
Since antiderivatives are so closely related to integrals by the
 Fundamental Theorems, we adopt the integral sign as a notation for the most general antiderivative of a function:
$$
\int f(x)\,dx \ = \ F(x)+C \quad\text{for all $C$.}
$$
Here $F(x)$ is a particular antiderivative: $F'(x)=f(x)$;
and $F(x)+C$ means the family of all antiderivatives, one
for every constant $C$ (\S3.9). This family is called the
{\it indefinite integral}, with no specific limits of integration
next to the $\int$ sign.
\\[.5em]
{\sc example:} Since $\frac{d}{dx}(x^3)=3x^2$ 
and $\frac{d}{dx}(\frac13x^3)=x^2$, we have the indefinite integral:
$$
\int x^2\,dx \ =\ \tfrac13x^3 + C\,.
$$
{\sc example:} Suppose a car with position function $s(t)$ 
lurches forward with velocity $v(t) =10t + 10\sin(\pi t)$ m/sec. 
How far does it travel ($\Delta s$) from $t=0$ to $t=3$ sec?
Using the antiderivative table in \S3.9, we first find
the indefinite integral:
$$
\int 10t+ 10\sin(\pi t)\,dt \ = \ 5t^2 -\tfrac{10}{\pi}\cos(\pi t)+C.
$$
Velocity is the rate of change of position, so total change in position is
the integral:
$$
\Delta s\ =\ s(3)-s(0) \ =\ \int_0^3 10t+ 10\sin(\pi t)\,dt
 \ = \ \left.5t^2 -\tfrac{10}{\pi}\cos(\pi t)\,\right|_{t=0}^{t=3}
$$
$$
\ =\ \left(5(3^2) -\tfrac{10}{\pi}\cos(\pi{\cdot} 3)\right)
- \left(5(0^2) -\tfrac{10}{\pi}\cos(\pi{\cdot} 0)\right)
\ =\ 45 + \tfrac{20}{\pi}
\ \approx\  51.4 \text{ meters.}
$$
\\[0em] 
{\bf Average of a function.} We can rewrite a Riemann sum with $\Delta x = \frac{b-a}n$  as:
$$
\sum_{i=1}^n f(x_i)\,\Delta x
\ = \ 
 \sum_{i=1}^n f(x_i)\,\frac{b{-}a}{n}
\ = \ 
(b{-}a)\,\frac{f(x_1)+\cdots +f(x_n)}{n}\,.
$$
This is just the interval length $(b{-}a)$
times the {\it average} of sample values $f(x_1),\ldots,f(x_n)$.
The integral is the limit of this as $n\to\infty$, which
becomes $(b{-}a)$ times the average of 
a denser and denser set of sample values:
$$
\int_a^b f(x)\,dx
\ =\ 
(b{-}a)\, \lim_{n\to\infty} \frac{f(x_1)+\cdots+f(x_n)}{n}\,.
$$

\addtolength{\textheight}{4em}

\pagebreak 

\addtolength{\topmargin}{-1em}

\noindent
We define the limit at right to be the continuous average of $f(x)$ over all $x\in[a,b]$:
$$
f_{\text{avg}} 
\ \ \stackrel{\text{def}}=\ \ 
\lim_{n\to\infty} \frac{f(x_1)+\cdots+f(x_n)}{n}
\ =\ 
\frac{1}{b{-}a}\,\int_a^b f(x)\,dx\,.
$$


\noindent
{\sc example:} Find the average value of $f(x) = \sqrt{x}$ over 
$x\in[0,4]$. The indefinite integral is: $\int \sqrt{x}\,dx = \int x^{1/2}\,dx
= \frac23 x^{3/2}+C$. The average is thus:
$$
\frac{1}{4{-}0}\,\int_0^4 \sqrt{x}\,dx
\ =\ \frac{1}{4}\,\left.\frac23 x^{\frac32}\,\right|_{x=0}^{x=4}
\ =\ \frac{4}{3}
\ \approx\ 1.3\ .
$$
Over $x\in[0,4]$ the function varies between $0\leq f(x) \leq 2$, 
but $f_{\text{avg}}\approx 1.3$  is above the midpoint 1.0,
because the graph bulges above the straight line from $(0,0)$ to $(4,2)$.
\\[1em]
\centerline{
\includegraphics[width=2.2in]
{4-4-Average.png}
}
\\[.3em]
A geometric way to picture the average is to think of the area under the curve as a fluid.
If we remove the curve $y=\sqrt x$ and contain the fluid between walls $x=0$ and $x=4$, 
then it drops to a horizontal level which is the average.
In the picture, the fluid under the straight line would
fill the container to the midpoint level $y=1$; but with 
extra fluid under $y=\sqrt x$ and above the line,  the average is higher: $y=\frac43$\,.
\\[.7em]
{\small 
{\sc example:} In  \S1.4, we intuitively took average velocity as distance traveled per time elapsed,
then defined instantaneous velocity as the short-time limit of this, leading to the definition of derivative. 
To check the validity of our definition, we reverse this reasoning:
starting with position and velocity functions $s(t)$ and $v(t)=s'(t)$, 
we compute the average velocity using the definition of average and the Second Fundamental Theorem:
$$\textstyle
v_{\text{avg}} \ =\  \frac{1}{b-a} \int_a^b v(t)\, dt 
 \ =\ \frac{1}{b-a} (s(b)-s(a))
 \ =\ \frac{s(b)-s(a)}{b-a}.
$$
That is, the average of $v(t)$ is indeed
distance traveled $s(b)-s(a)$ divided by the time elapsed $b-a$.
This expected result confirms that our axiomatic theory captures
the physical intuition that motivated it, making us more confident in 
applying the theory to make physical predictions beyond our intuition.
}
\\[1em]
{\small 
{\bf Mean Value Theorem for Integrals} 
\begin{quote}
If $f(x)$ is continuous for $x\in[a,b]$,
then there is some $c\in(a,b)$ where $f(c)$ equals 
the average of $f(x)$ over the interval:
$f(c)=  f_{\text{avg}}=$~$\frac{1}{b{-}a}\,\int_a^b f(x)\,dx$.
\end{quote}

\noindent
{\it Proof:} Take $I(x) =\int_a^x f(x)\,dx$. The Mean Value Theorem for Derivatives 
(\S3.2) says there is a value $c\in(a,b)$ where the tangent line for $I(x)$ 
is parallel to the secant line over the interval: $I'(c) = \frac{I(b)-I(a)}{b-a}$. By the First Fundamental Theorem, the 
left side is $I'(c) = f(c)$; and since $I(a) = 0$, the right side is $\frac{I(b)}{b-a} =\frac{1}{b-a}\int_a^b f(x)\,dx$, as desired.
\\[0em]

\noindent
In our example $f(x) = \sqrt{x}$, there is $c\in(0,4)$ with 
$f(c)=\sqrt c = f_{\text{avg}}=\frac43$: i.e.~$c=\frac{16}{9}$.
}





\end{document}

%%%%%%%%%%%  Picture  %%%%%%%%%%
\\[1em]
\centerline{
%\includegraphics[width=2in]
{1-8-Discont-iv.png}
\qquad \qquad 
%\includegraphics[width=2in]
{1-8-Discont-v.png} 
}
\vspace{1em}

\noindent

%%%%%%%%%%%  Table  %%%%%%%%%%
$$\begin{array}{|c|c|}
A & B
\\ \hline && \\[-1.1em]
C & D
\end{array}$$


{\bf Review.} The integral of $y=f(x)$ has four levels of meaning:
\begin{itemize}
\item Physical: If $y=\frac{dz}{dx}$ is the rate of change of some other quantity $z$, 
then $\int_a^b f(x)\,dx$ is the total change in $z$ 
from $x=a$ to $x=b$.

For example, if $v=\frac{ds}{dt}$ is a velocity, then the integral is the total change in position $s$, 
which is the distance traveled: $\int_a^b v(t)\,dt = s(b)-s(a)$.

\item Geometric: $\int_a^b f(x)\,dx$ is the area between the curve $y=f(x)$ and the
interval $[a,b]$ on the $x$-axis, provided $f(x)\geq 0$.\footnote{
We shall see that this remains valid for $f(x)\leq 0$ if we count area below the $x$-axis as {\it negative}.}

\item Numerical: To define the integral as a limit of approximations, we split the interval $x\in[a,b]$ into $n$ increments of size $\Delta x=\frac{b-a}{n}$, and choose a sample point in each increment:
$x_1,x_2,\ldots, x_n$ can be the left or right endpoints, or anywhere between. Then:
$$
\int_a^b f(x)\,dx\ \ =\ \lim_{\Delta x\to 0} f(x_1)\Delta x + f(x_2)\Delta x+\cdots +f(x_n)\Delta x.
$$
{\it First Fundamental Theorem:} If we fix $a$, and we define $F(x)$ for $x=b\geq a$ by $F(b)\ =\ \int_a^b f(x)\,dx$, then $F(x)$ is
an antiderivative of $f(x)$ with $F(a)=0$.\footnote{
We shall see that this remains valid for $b\leq a$ if we define $\int_a^b f(x)\,dx$ as $-\int_b^a f(x)\,dx$.
}
\item Algebraic: If we can reverse differentiation rules to get $f(x)$ as the derivative of a known function $F(x)$,
then $\int_a^b f(x)\,dx=F(b)-F(a)$, the change in $F(x)$ from $x=a$ to $x=b$.
This is known as  the {\it Second Fundamental Theorem}.
\end{itemize}

\begin{quote}
{\it First Fundamental Theorem of Calculus.} Given a continuous function $f(x)$, 
fix some $x=a$, and define $F(x)$ for any $x=b\geq a$ as:
$$
F(b) \ =\ \int_a^b\! f(x) \,dx\,.
$$
Then $F'(x)=f(x)$, and $F(x)$ is an antiderivative of $f(x)$ with initial value $F(a)=0$.
\end{quote}
%
We will discuss this Theorem in detail later, but it should be clear that we defined
the integral just so as to make it true. 
\\[1em]




